% !TEX root = ../main.tex
% File: chapters_part1/chap7_3.tex
% Nội dung cho Chương 7, Phần 3

\section{Thách thức trong đánh giá LLM}
\label{sec:llm_evaluation_challenges}

Đánh giá mô hình ngôn ngữ lớn vẫn là một bài toán mở, dù bộ metric và benchmark ngày càng tinh vi. Một con số duy nhất trên bảng xếp hạng (leaderboard) có thể không phản ánh hiệu năng thực sự của mô hình khi triển khai.

\subsection{Nhiễm bẩn dữ liệu (Data Contamination)}
\label{ssec:data_contamination}

\begin{pitfall}{Vấn đề cốt lõi}
Điểm số trên một benchmark bị nhiễm bẩn không cho biết mô hình đang suy luận hay chỉ đang nhớ lại đáp án đã gặp khi huấn luyện.
\end{pitfall}

\paragraph{Định nghĩa}
Nhiễm bẩn dữ liệu xảy ra khi dữ liệu của bộ \textbf{kiểm tra (test set)} hoặc bộ \textbf{phát triển (dev set)} của một benchmark vô tình xuất hiện trong bộ \textbf{dữ liệu huấn luyện (training set)} của mô hình.

\paragraph{Nguyên nhân}
\begin{itemize}
    \item \textbf{Dữ liệu pre-training quá lớn:} Các LLM được huấn luyện trên một phần lớn của Internet, ví dụ Common Crawl. Các bộ dữ liệu benchmark phổ biến vốn cũng thường được đăng công khai trên GitHub hay arXiv, nên rất có khả năng chúng đã nằm sẵn trong kho dữ liệu đó.
    \item \textbf{Trùng lặp không chủ ý:} Ngay cả khi không trùng lặp trực tiếp, các câu hỏi và câu trả lời trong benchmark có thể có phiên bản diễn giải lại (paraphrased versions), hoặc các thảo luận về chúng trên Stack Overflow hay Reddit, và mô hình đã học được chúng.
\end{itemize}

\paragraph{Hậu quả}
\begin{itemize}
    \item \textbf{Điểm số bị thổi phồng:} Khi gặp một câu hỏi đã “nhìn thấy” trong quá trình huấn luyện, mô hình không cần suy luận mà chỉ cần nhớ lại và đưa ra câu trả lời. Điểm số vì thế cao hơn năng lực thực.
    \item \textbf{Benchmark mất giá trị:} Một benchmark bị nhiễm bẩn nặng không còn đo được khả năng tổng quát hóa.
\end{itemize}

\paragraph{Hướng giải quyết}
Chưa có cách giải quyết triệt để. Các hướng tiếp cận gồm:
\begin{itemize}
    \item \textbf{Lọc dữ liệu:} Loại bỏ dữ liệu của các benchmark đã biết khỏi tập pre-training.
    \item \textbf{Liên tục phát triển benchmark mới:} Tạo các bộ dữ liệu mới và giữ kín cho đến khi công bố.
    \item \textbf{Phát hiện nhiễm bẩn:} Xây dựng các kỹ thuật ước tính mức độ một mô hình đã học thuộc lòng một benchmark.
\end{itemize}

\subsection{Đánh giá bởi con người so với đánh giá tự động (Human vs. Auto Evaluation)}
\label{ssec:human_vs_auto_eval}

\subsubsection{Hạn chế của đánh giá tự động (Auto-Eval)}
Như đã thảo luận ở mục \ref{sec:classic_metrics}, các metric tự động như BLEU, ROUGE chỉ dựa trên sự trùng lặp bề mặt của từ hoặc N-gram. Chúng không đánh giá được sự sáng tạo, tính mạch lạc hay độ chính xác thông tin: một câu trả lời sai hoàn toàn về mặt thông tin vẫn có thể đạt điểm ROUGE cao nếu nó dùng lại các từ trong tài liệu tham khảo. Chúng cũng không đo được “tính hữu ích”, “sự thú vị” hay “mức độ an toàn” của một câu trả lời.

\subsubsection{Ưu điểm và thách thức của đánh giá bởi con người (Human-Eval)}
Người đánh giá nắm bắt được những khía cạnh tinh tế mà metric tự động bỏ qua, nên đánh giá của con người được dùng làm mốc tham chiếu. Cái giá là ba vấn đề. Thứ nhất, nó tốn kém và chậm, vì cần thuê và huấn luyện một đội ngũ người đánh giá chuyên nghiệp. Thứ hai, nó thiếu nhất quán: những người đánh giá khác nhau có sở thích và tiêu chuẩn khác nhau, dẫn tới mức đồng thuận thấp giữa những người gán nhãn (low inter-annotator agreement). Thứ ba, nó có thiên vị (bias): người đánh giá bị ảnh hưởng bởi độ dài câu trả lời, thường ưu ái câu dài hơn, hoặc bởi văn phong của mô hình.

Cách làm thông dụng là kết hợp cả hai: dùng metric tự động để theo dõi nhanh trong quá trình phát triển, và dùng đánh giá của con người theo định kỳ ở các mốc quan trọng.

\subsection{Định luật Goodhart và “dạy cho bài kiểm tra” (Goodhart's Law \& ``Teaching to the Test'')}
\label{ssec:goodharts_law}

\begin{keyidea}{Định luật Goodhart}
“When a measure becomes a target, it ceases to be a good measure.” \\ (Khi một thước đo trở thành một mục tiêu, nó sẽ không còn là một thước đo tốt nữa.)
\end{keyidea}

\paragraph{Ý nghĩa trong bối cảnh LLM}
Khi một benchmark cụ thể, chẳng hạn MMLU, trở thành mục tiêu chính của các phòng thí nghiệm AI, các nhà phát triển bắt đầu tối ưu mô hình trực tiếp cho benchmark đó (over-optimizing). Họ có thể vô tình hoặc cố ý đưa dữ liệu có văn phong tương tự benchmark vào quá trình huấn luyện, hoặc tinh chỉnh siêu tham số để mô hình hoạt động tốt nhất trên phân phối dữ liệu rất hẹp của benchmark. Kết quả là điểm số trên benchmark tăng cao mà khả năng suy luận tổng quát trong các tình huống thực tế không tăng theo: mô hình chỉ học được cách “làm bài kiểm tra”.

Hệ quả là các bảng xếp hạng dần không còn phản ánh đúng tiến bộ của ngành, và các phương pháp đánh giá phải liên tục được làm mới và đa dạng hóa.